\usepackage{float}
% preamble.tex -- a PE nyílt tankönyvsorozat PDF-arculata (lualatex, scrbook)
% Betűk: IBM Plex Serif (törzsszöveg), IBM Plex Sans (feliratok), Raleway (címek).

% --- Betűk --------------------------------------------------------------
\setmainfont{IBM Plex Serif}[BoldFont={IBM Plex Serif SemiBold},
  ItalicFont={IBM Plex Serif Italic}, BoldItalicFont={IBM Plex Serif SemiBold Italic}]
\setsansfont{IBM Plex Sans}[BoldFont={IBM Plex Sans SemiBold},
  ItalicFont={IBM Plex Sans Italic}, BoldItalicFont={IBM Plex Sans Bold Italic}]
\setmonofont{IBM Plex Mono}[Scale=0.88]
\newfontfamily\ralewayfekete{Raleway Black}[Ligatures=TeX]
\newfontfamily\ralewayextra{Raleway ExtraBold}[BoldFont={Raleway Black}, Ligatures=TeX]
\newfontfamily\ralewaybold{Raleway Bold}[Ligatures=TeX]
\newfontfamily\ralewaynormal{Raleway}[BoldFont={Raleway Bold}, Ligatures=TeX]

% --- Színek (azonosak a peabra.sty palettájával) -------------------------
\usepackage{xcolor}
\definecolor{peKek}{HTML}{285780}
\definecolor{peSotet}{HTML}{1B3A57}
\definecolor{peKozep}{HTML}{4F7FAD}
\definecolor{peVilagos}{HTML}{A9C4DE}
\definecolor{peHalvany}{HTML}{E6EEF6}
\definecolor{peKiemel}{HTML}{E07B39}
\definecolor{peKiemelH}{HTML}{FBE7D9}
\definecolor{peZold}{HTML}{3E8E63}
\definecolor{peZoldH}{HTML}{DDEFE4}
\definecolor{pePiros}{HTML}{B83A3A}
\definecolor{pePirosH}{HTML}{F6DEDE}
\definecolor{peLila}{HTML}{7A5C99}
\definecolor{peLilaH}{HTML}{ECE5F3}
\definecolor{peSarga}{HTML}{D9A520}
\definecolor{peSzurke}{HTML}{6B7280}
\definecolor{peSzurkeH}{HTML}{F1F2F4}
\definecolor{peSzoveg}{HTML}{1F2933}

% --- KOMA-Script címsorok és élőfej --------------------------------------
\setkomafont{disposition}{\ralewayextra\color{peSotet}}
\setkomafont{part}{\ralewayfekete\Huge\color{peSotet}}
\setkomafont{partnumber}{\ralewayextra\LARGE\color{peKiemel}}
\setkomafont{chapter}{\ralewayfekete\huge\color{peSotet}}
\setkomafont{section}{\ralewayextra\Large\color{peSotet}}
\setkomafont{subsection}{\ralewayextra\large\color{peKek}}
\setkomafont{subsubsection}{\ralewaybold\normalsize\color{peKek}}
\setkomafont{paragraph}{\ralewaybold\normalsize\color{peSotet}}
\renewcommand*{\chapterformat}{\mbox{\color{peKiemel}\thechapter\autodot}\enskip}
\setkomafont{captionlabel}{\sffamily\bfseries\color{peKek}}
\setkomafont{caption}{\sffamily\small}
\setkomafont{pageheadfoot}{\sffamily\footnotesize\color{peSzurke}}
\setkomafont{pagenumber}{\sffamily\footnotesize\color{peSotet}}
\RedeclareSectionCommand[beforeskip=-1.2\baselineskip, afterskip=0.5\baselineskip]{section}

% --- Kiemelődobozok (a dobozok.lua szűrő hívja) --------------------------
\usepackage[most]{tcolorbox}
\colorlet{dobozF-cel}{peKek}         \colorlet{dobozH-cel}{peKek!7}
\colorlet{dobozF-fogalom}{peSotet}   \colorlet{dobozH-fogalom}{peSotet!6}
\colorlet{dobozF-eset}{peKozep}      \colorlet{dobozH-eset}{peKozep!8}
\colorlet{dobozF-esetA}{peKek}       \colorlet{dobozH-esetA}{peKek!7}
\colorlet{dobozF-esetB}{peZold}      \colorlet{dobozH-esetB}{peZold!9}
\colorlet{dobozF-esetC}{peLila}      \colorlet{dobozH-esetC}{peLila!9}
\colorlet{dobozF-pelda}{peKozep}     \colorlet{dobozH-pelda}{peKozep!6}
\colorlet{dobozF-kitekintes}{peLila} \colorlet{dobozH-kitekintes}{peLila!7}
\colorlet{dobozF-hiba}{pePiros}      \colorlet{dobozH-hiba}{pePiros!7}
\colorlet{dobozF-onell}{peZold}      \colorlet{dobozH-onell}{peZold!7}
\colorlet{dobozF-projekt}{peKiemel}  \colorlet{dobozH-projekt}{peKiemel!9}
\colorlet{dobozF-prompt}{peSzurke}   \colorlet{dobozH-prompt}{peSzurke!7}
\colorlet{dobozF-osszefoglalo}{peKek}\colorlet{dobozH-osszefoglalo}{peKek!7}
\colorlet{dobozF-forras}{peSzurke}   \colorlet{dobozH-forras}{peSzurke!5}
\colorlet{dobozF-labor}{peKiemel}    \colorlet{dobozH-labor}{peKiemel!8}
\colorlet{dobozF-baleset}{pePiros}   \colorlet{dobozH-baleset}{pePiros!5}
\newtcolorbox{pedoboz}[2]{enhanced, breakable, sharp corners,
  boxrule=0pt, leftrule=2.6pt, toprule=0pt, rightrule=0pt, bottomrule=0pt,
  colframe=dobozF-#1, colback=dobozH-#1, colbacktitle=dobozH-#1, coltitle=dobozF-#1,
  title={#2}, fonttitle=\ralewayextra\small, titlerule=0pt,
  toptitle=1.6mm, bottomtitle=-0.4mm, left=3.2mm, right=3.2mm, top=1mm, bottom=1.8mm,
  before skip=3.5mm, after skip=3.5mm, fontupper=\small}
\newenvironment{pemegoldas}{\par\smallskip\noindent\begingroup\footnotesize
  {\sffamily\bfseries\color{peZold}Megoldás.}\ }{\par\endgroup}

% --- Egyéb ---------------------------------------------------------------
\usepackage{tikz}
\usetikzlibrary{calc}
\usepackage{booktabs}
\AtBeginDocument{\hypersetup{linkcolor=peKek, urlcolor=peKek, citecolor=peKek}}
\setlength{\emergencystretch}{3em}
\widowpenalty=10000
\clubpenalty=10000

% --- Arculatos címoldal ferde sávval -------------------------------------
\makeatletter
\renewcommand{\maketitle}{%
  \begin{titlepage}
  \thispagestyle{empty}
  \begin{tikzpicture}[remember picture, overlay]
    \fill[peKek] (current page.south west) -- (current page.south east)
      -- ($(current page.south east)+(0,10.5cm)$) -- ($(current page.south west)+(0,7cm)$) -- cycle;
    \fill[peKiemel] ($(current page.south west)+(0,7.35cm)$) -- ($(current page.south east)+(0,10.85cm)$)
      -- ($(current page.south east)+(0,11.15cm)$) -- ($(current page.south west)+(0,7.65cm)$) -- cycle;
    \node[anchor=north west, text width=14cm, align=left, inner sep=0]
      at ($(current page.north west)+(2.4cm,-5.2cm)$) {%
      {\ralewayfekete\fontsize{40}{44}\selectfont\color{peSotet}\@title\par}
      \vspace{7mm}
      {\ralewayextra\fontsize{17}{22}\selectfont\color{peKek}\@subtitle\par}
      \vspace{16mm}
      {\ralewaynormal\Large\color{peSzoveg}\@author\par}};
    \node[anchor=south west, text width=14cm, align=left, inner sep=0, text=white]
      at ($(current page.south west)+(2.4cm,2.6cm)$) {%
      {\ralewayextra\large Pannon Egyetem\par}\vspace{1mm}{\sffamily\small Veszprém, \@date\par}};
  \end{tikzpicture}
  \end{titlepage}%
}
\makeatother

% A dobozokon belül az ábrák nem úszhatnak: helyben, minipage-ben jelennek meg.
\makeatletter
\newcommand{\pefigurahelyben}{%
  \renewenvironment{figure}[1][]{\par\medskip\noindent\begin{minipage}{\linewidth}\def\@captype{figure}\centering}%
    {\end{minipage}\par\medskip}}
\makeatother

% --- Magyar címkék, élőfej és ábrafeliratok -------------------------------
% (a babel magyar nyelvi fájlja nélkül is helyes magyar formák)
\renewcommand*{\chaptermarkformat}{\thechapter\autodot\enskip}
\renewcommand*{\partformat}{}
\renewcommand*{\addparttocentry}[2]{\addtocentrydefault{part}{}{#2}}
\AtBeginDocument{%
  \renewcommand{\chaptername}{fejezet}\renewcommand{\partname}{rész}%
  \renewcommand{\appendixname}{Függelék}\renewcommand{\contentsname}{Tartalom}%
  \renewcommand{\bibname}{Irodalom}\renewcommand{\figurename}{ábra}\renewcommand{\tablename}{táblázat}%
  \DeclareCaptionLabelFormat{pehu}{#2.~#1}%
  \captionsetup{labelformat=pehu,labelsep=period}}
% Görög betűk a folyó szövegben (a törzsbetűben nincs görög jelkészlet)
\usepackage{newunicodechar}
\newunicodechar{α}{\ensuremath{\alpha}}\newunicodechar{β}{\ensuremath{\beta}}
\newunicodechar{γ}{\ensuremath{\gamma}}\newunicodechar{δ}{\ensuremath{\delta}}
\newunicodechar{ε}{\ensuremath{\varepsilon}}\newunicodechar{η}{\ensuremath{\eta}}
\newunicodechar{λ}{\ensuremath{\lambda}}\newunicodechar{μ}{\ensuremath{\mu}}
\newunicodechar{π}{\ensuremath{\pi}}\newunicodechar{σ}{\ensuremath{\sigma}}
\newunicodechar{χ}{\ensuremath{\chi}}\newunicodechar{Γ}{\ensuremath{\Gamma}}
\newunicodechar{Δ}{\ensuremath{\Delta}}\newunicodechar{Λ}{\ensuremath{\Lambda}}
\newunicodechar{Σ}{\ensuremath{\Sigma}}\newunicodechar{Ω}{\ensuremath{\Omega}}
\raggedbottom
